% CV · Oumayma El Moussaoui
% Compile with: pdflatex cv.tex
% The phone number lives in personal.tex, which is kept out of the public repo.
% Without it the CV still compiles (the phone number is left out).
\documentclass[10pt,a4paper]{article}

\usepackage[a4paper,top=11mm,bottom=11mm,left=14mm,right=14mm]{geometry}
\input{glyphtounicode}\pdfgentounicode=1   % correct Unicode text for ATS parsers
\pdfinterwordspaceon \pdffakespace          % real space characters between words
\usepackage[T1]{fontenc}
\usepackage[utf8]{inputenc}
\usepackage{lmodern}
\usepackage{microtype}
\usepackage{titlesec}
\usepackage{enumitem}
\usepackage[hidelinks]{hyperref}
\hypersetup{pdftitle={CV - Oumayma El Moussaoui},pdfauthor={Oumayma El Moussaoui},
  pdfsubject={Curriculum Vitae},pdfkeywords={Data Scientist, Machine Learning, Computer Vision, Python}}

\pagestyle{empty}
\setlength{\parindent}{0pt}
\setlength{\parskip}{0pt}
\raggedright
\hyphenpenalty=10000
\exhyphenpenalty=10000
\emergencystretch=3em

% Section headings: large, with a full-width rule underneath
\titleformat{\section}{\Large\bfseries}{}{0pt}{}[\vspace{1pt}\titlerule]
\titlespacing*{\section}{0pt}{8pt}{3pt}

% Bullets
\setlist[itemize]{leftmargin=12pt,labelsep=5pt,itemsep=0.5pt,topsep=1pt,parsep=0pt,partopsep=0pt}

% \org{Organisation}{Dates}{Role or unit}{Location}
\newcommand{\org}[4]{%
  \par\vspace{4pt}\noindent\textbf{#1}\hfill #2\par
  \noindent\textit{#3}\hfill\textit{#4}\par}
% \role{Role}{Dates}: a position within the organisation above
\newcommand{\role}[2]{\par\vspace{1pt}\noindent\textbf{#1}\hfill\textit{#2}\par}
% \project{Name}{Tech stack}
\newcommand{\project}[2]{\par\vspace{4pt}\noindent\textbf{#1}\hfill\textit{#2}\par}

\IfFileExists{personal.tex}{\input{personal.tex}}{\newcommand{\cvphone}{}}
\newcommand{\sep}{\ \textbar\ }

\begin{document}

% ---------- header ----------
\begin{center}
  {\fontsize{24}{28}\selectfont Oumayma El Moussaoui}\par\vspace{4pt}
  {\small
  Aalen, Germany\sep \ifx\cvphone\empty\else\cvphone\sep\fi
  \href{mailto:oma98yma@gmail.com}{oma98yma@gmail.com}\sep
  \href{https://www.linkedin.com/in/oumayma-el-9aa463375}{linkedin.com/in/oumayma-el-9aa463375}\par
  \href{https://github.com/elmoussaouioumayma}{github.com/elmoussaouioumayma}\sep
  \href{https://elmoussaouioumayma.github.io/Portfolio/}{elmoussaouioumayma.github.io/Portfolio}\par}
  \vspace{4pt}
  \textbf{Data Scientist \textbar\ Image \& Measurement Data \textbar\ Engineering Background}\par
\end{center}
\vspace{-6pt}

% ---------- profile ----------
\section{Profile}
Data scientist with two master's degrees, one in \textbf{Machine Learning \& Data Analytics} (Hochschule Aalen) and one in \textbf{Energy Engineering}, and industrial experience in cement, mining and manufacturing plants. I turn raw image and measurement data into \textbf{reliable numbers that engineers can act on}. I work in \textbf{Python}, validate my results carefully, and am used to working with chemists, materials scientists and process engineers.

% ---------- experience ----------
\section{Experience}
\org{Hochschule Aalen}{Nov 2024 -- present}{Zentrum Elektrochemische Oberflächentechnik (ZEO)}{Aalen}
\role{Academic Employee (Wissenschaftliche Mitarbeiterin)}{Feb 2026 -- present}
\role{Student Research Assistant}{Nov 2024 -- Jan 2026}
\begin{itemize}
  \item Develop \textbf{image-analysis software} (Python, OpenCV, scikit-image) that automatically \textbf{detects, classifies and measures cracks} in microscope images of chromium coatings.
  \item Build \textbf{interactive tools} that researchers can use themselves: live parameter tuning, batch processing with CSV export, depth and width profiles along cracks.
  \item Analyse process and measurement data from chromium electroplating experiments and \textbf{present the results to chemists and engineers}.
  \item Model electroplating process knowledge as an \textbf{ontology (Protégé/OWL)} so that experimental data can be structured and reused.
\end{itemize}

\org{Excellence Hub}{Sep 2021 -- Nov 2021}{Engineering Intern: Energy Efficiency}{El Jadida, Morocco}
\begin{itemize}
  \item Carried out an \textbf{energy audit} of an industrial plant (MAZAFIL): collected and analysed consumption data, \textbf{quantified the energy losses} and proposed savings measures.
\end{itemize}

\org{LafargeHolcim}{Apr 2019 -- Jun 2019}{Engineering Intern: Thermal Processes}{Meknes, Morocco}
\begin{itemize}
  \item Analysed the \textbf{thermal losses of a cement kiln} and evaluated insulation options to reduce them.
\end{itemize}

\org{OCP Group}{Apr 2018 -- May 2018}{Technical Intern: Process Performance}{Khouribga, Morocco}
\begin{itemize}
  \item Monitored the operation of an industrial grinder and proposed \textbf{cost-effective measures} to improve its performance.
\end{itemize}

% ---------- projects ----------
\section{Projects (github.com/elmoussaouioumayma)}
\project{Master's Thesis: Quantitative Crack Analysis in Chromium Layers}{Python, OpenCV, scikit-image}
\begin{itemize}
  \item \textit{Optimized Image Processing for Quantitative Crack Analysis in Chromium Layers} (Hochschule Aalen, 2026): automated Python pipeline that turns microscope images of hard-chromium coatings into crack measurements, \textbf{replacing manual annotation (2 to 2.5 hours per image)}.
  \item Segmentation (5 classical methods compared, top-hat and Otsu selected), \textbf{vectorization into a crack graph}, then crack length, density and junction angles in µm.
  \item Validated against manual ground truth: \textbf{chamfer distance under 1.77 px (approx. 0.08 µm)} on 7 of 8 images. Crack density on 20 images (564 +/- 70 cracks/cm) lies within the published range for hard chromium.
\end{itemize}

\project{Crack Depth Explorer: 3D Crack Measurement from Four Photos}{Python, NumPy/SciPy, OpenCV}
\begin{itemize}
  \item Measures cracks in 3D (depth, width and length in mm) from four photos taken with a fixed camera and moving light: \textbf{photometric stereo} reconstructs the surface, then every crack is detected and measured. For concrete, stone, plaster and asphalt; runs live in the browser.
  \item On 20 test scenes with known depth: \textbf{0.21 mm mean depth error}, \textbf{2.6 times more accurate} than single-photo methods, detection precision / recall 0.96 / 0.97.
\end{itemize}

\project{Crack Network Graph: From Crack Images to Measured Graphs}{Python, OpenCV, scikit-image, NetworkX}
\begin{itemize}
  \item Turns a photo of a cracked surface into a \textbf{measured graph} (pixels to vectors): crack lengths, widths and directions, T / Y / X junctions and the closed cells between cracks; the junction types show how the network formed. Exports GeoJSON and GraphML; runs live in the browser.
  \item On 30 images with a known network: \textbf{90\% of junction types correct}, formation type recognised in \textbf{30 of 30 images}.
\end{itemize}

\project{Coating Layer Analysis}{Python, image analysis, statistical modelling}
\begin{itemize}
  \item Detects the layers of a coating in cross-section images from changes in particle size and density, traces the wavy boundaries and measures the thickness of each layer; the number of layers is chosen with a \textbf{statistical model (BIC)}. Includes a web app and a batch tool.
  \item On 100 test images with known layers: \textbf{95\% correct number of layers}, median boundary error 0.66 µm.
\end{itemize}

\project{Dark Pattern Detection in User Interfaces}{Python, NLP, OCR, Gradio}
\begin{itemize}
  \item Classifier for manipulative website texts. \textbf{Cross-validation grouped by website} prevents leakage between training and test data (\textbf{test F1 97.3\%}). Gradio app that also reads screenshots via OCR.
\end{itemize}

\project{Fake vs. Real News Detection (NLP)}{Python, scikit-learn, TF-IDF}
\begin{itemize}
  \item \textbf{Found and removed data leakage} (5,498 duplicates and source markers that revealed the label), then built a TF-IDF and Logistic Regression model with \textbf{97.1\% test accuracy} and tested it on newly scraped articles.
\end{itemize}

\project{Diabetes Prediction}{Python, scikit-learn, XGBoost}
\begin{itemize}
  \item Compared 6 models in \textbf{leak-free pipelines (CV AUC 0.84)} and tuned the decision threshold for screening, reaching \textbf{85\% recall} on the test set.
\end{itemize}

% ---------- skills ----------
\section{Skills}
\textbf{Programming:} Python (pandas, NumPy, SciPy), SQL, R, Git/GitHub, Jupyter\par
\textbf{Image processing:} OpenCV, scikit-image: segmentation, morphological filtering, skeletonization, graph extraction, geometric measurement\par
\textbf{Machine learning:} scikit-learn, XGBoost, NLP (TF-IDF, Hugging Face Transformers); cross-validation, leakage checks, threshold tuning\par
\textbf{Data \& visualisation:} Power BI, matplotlib, interactive tools (Gradio, matplotlib widgets), technical reporting\par
\textbf{Engineering:} energy audits, thermal analysis, COMSOL, PVsyst, AutoCAD, SolidWorks; ontologies (Protégé/OWL)\par
\textbf{Languages:} Arabic (native), French (fluent), English (fluent), German (B2, telc certificate)\par

% ---------- education ----------
\section{Education}
\org{Hochschule Aalen}{2024 -- 2026}{M.Sc. Machine Learning and Data Analytics}{Aalen, Germany}
\org{Intensive German Language Course}{2022 -- 2024}{telc Deutsch B2; preparation for master's studies in Germany}{}
\org{ENSIAS}{2021 -- 2022}{Doctoral Studies in Energy Engineering (1 year of research, not continued)}{Rabat, Morocco}
\noindent Research on radiative sky cooling with emissive materials.\par
\org{ENSAM}{2019 -- 2021}{Master's in Energy and Environmental Sciences}{Rabat, Morocco}
\noindent Thesis: fabrication and characterisation of emissive materials and testing of a night-time radiative cooling test bench.\par
\org{ENS}{2018 -- 2019}{Bachelor's in Physical Energy Technologies}{Rabat, Morocco}
\org{Technician Diploma in Energy Technologies}{2016 -- 2018}{}{Errachidia, Morocco}

\vspace{5pt}
\textbf{Certificates:} Data Engineer, DataCamp (2023); Data Analyst, Google (2022); First Aid, Red Crescent (2023).\par
\textbf{Volunteering:} Finance manager of a cultural and social committee (2021); ALBASMA medical and humanitarian caravan (2020).\par

\end{document}
